\section{\lang: Guardrail Description Language}\label{sec:gdl}

We present the \emph{Guardrails Description Language} (GDL), a declarative language for specifying performance guardrails over kernel behavior.
As discussed in \Cref{subsec:semantics}, a guardrail is given by: (i) \emph{what condition} must hold on the system state it inspects (the monitorable property $\property{G}$), (ii) \emph{what corrective action} is applied when that condition fails ($\action{G}$), and (iii) \emph{which policies} the guardrail applies to ($\policies_G$).
Alongwith the guardrail specification, the system also needs to know the local guardrail properties $\varphi$ and the global \perfsafety property $\Phi$ that the guardrail is designed to enforce.
Thus, guardrail developers (OS operators or OS policy developers) are responsible for providing all these pieces of information.
In this section, we describe the primitives \lang introduces to make this specification intuitive.

A crucial aspect of guardrail specification is determining \emph{when} to evaluate the guardrail.
In principle, any point in the execution sequence of state transitions, decisions, and events can be used for evaluation.
However, there is a cost associated with evaluating the run-time performance property -- of reading the relevant state, computing the condition, and then, taking the action (if needed).
Therefore, the question is {\it when should the system incur the cost of evaluating a guardrail?}
In \lang, users can control this overhead by specifying explicit {\em triggers}, which determine \emph{when} the system should check the guardrail property.

Concretely, a guardrail in \lang is declared using six keywords corresponding: {\sf for}, {\sf when}, {\sf check}, {\sf action}, {\sf ensures}, and {\sf assert}.
The keyword {\sf for} names the set of policies for which the guardrail is relevant.
{\sf when} specifies the \emph{trigger}: the point in the execution at which the guardrail is evaluated.
{\sf check} encodes the performance condition.
{\sf action} prescribes the corrective action taken when the condition fails.
{\sf ensures} specifies the local guardrail properties $\varphi$ that must hold.
Finally, {\sf assert} specifies the global \perfsafety property $\Phi$ that the guardrail is designed to enforce.
\Cref{tab:gdl-primitives} summarizes the key primitives available in \lang, organized by their role.
Below, we describe these primitives in more detail and illustrate their use with examples of guardrails in \lang.

\begin{table}[t]
\centering
\small
\setlength{\tabcolsep}{4pt}
\begin{tabular}{p{1.4cm}p{2cm}p{4.4cm}}
\toprule
\textbf{Category} & \textbf{Primitive} & \textbf{Description} \\
\midrule
\multirow{2}{*}{Trigger}
  & $\textsf{periodic}(d)$ & Evaluate every $d$ seconds \\
  & $\textsf{on\_call}(f)$ & Evaluate on each invocation of \textsf{KFunc} $f$ \\
\midrule
\multirow{8}{1.4cm}{Inspection Objects}
  & $\textsf{KObject}$ & Kernel objects (pages, processes, \ldots) \\
  & $Set[\textsf{KObject}]$ & Collections of kernel objects \\
  & $\textsf{attribute}(x)$ & Attribute of a \textsf{KObject} or set \\
  & $\textsf{exec\_time}(f)$ & Execution time of \textsf{KFunc} $f$ \\
  & $\textsf{decision}(f)$ & Decision value returned by \textsf{KFunc} $f$ \\
  & $\textsf{arg}(f, \mathit{name})$ & Argument value at a \textsf{KFunc} $f$ invocation \\
  & $\textsf{read}(\mathit{feat})$ & Read a feature from the in-kernel store \\
  & \makecell[l]{$\textsf{avg}$/$\textsf{min}$/\\$\textsf{max}$/$\textsf{sum}(\cdot)$} & Aggregate an observation \\
\midrule
\multirow{2}{*}{Action}
  & $\textsf{Set}(\mathit{flag},\,\mathit{val})$ & Update a control flag or parameter \\
  & $\textsf{Report}$ & Log the violation for offline diagnosis \\
\bottomrule
\end{tabular}
\caption{Key primitives in the Guardrails Description Language (\lang).}
\label{tab:gdl-primitives}
\end{table}


\paragraph{Triggers.}
\lang supports two trigger types that cover the two most common classes of performance properties.
The first class of properties specifies continuous conditions that must always hold, such as CPU utilization by kernel daemons remaining below a threshold, or memory compaction overhead staying bounded.
Such properties are not tied to specific state transitions and are best checked at regular intervals.
\lang captures this with the {\sf periodic}$(d)$ trigger, which fires every $d$ seconds.
For example, the {\tt CompactionOverhead} guardrail from \Cref{subsec:compaction-case-study} uses a 10-second periodic trigger to check whether the average execution time of {\tt compact\_zone()} has grown excessive:

\begin{dslbox}
for THPAlways
when periodic(10s)
check $\forall e \in$ compact_zone() avg(exec_time(e)) < 10ms
action Set(thp_policy, THPNever)
\end{dslbox}

The second class of properties evaluates the correctness or efficacy of a particular policy \emph{decision}: e.g., whether a flow-size prediction used features that are in-distribution for the model, or whether a page promotion moved a cold page into the hot tier.
For such properties, the natural evaluation point is the moment the decision is produced---at the invocation of the relevant kernel function.
\lang captures this with the {\sf on\_call}$(f)$ trigger, which fires on each call to \textsf{KFunc} $f$.
For example, the {\tt InDistributionShort} guardrail from \Cref{subsec:flux-case-study} fires on each call to {\tt flow\_start} and checks whether the input features lie within the training distribution of the active predictor:

\begin{dslbox}
for FLUXShortPredictor
when on_call(flow_start)
check $agg\_net \sim$ read(ShortDist[agg_net]) $\wedge$
      $cpu \sim$ read(ShortDist[cpu]) $\wedge\ \dots$
action Set(flux_policy, FLUXLongPredictor)
\end{dslbox}

There may be properties tied to higher-level events, where the state transition of interest may be missed by periodic sampling and cannot be captured cleanly by any single function invocation.
Such cases can be handled by encoding the event as a predicate over observations and evaluating the guardrail periodically at a sufficiently fine granularity.
Concretely, if a guardrail is intended to enforce $\property{}$ each time an event $E$ occurs, one can equivalently specify it with a periodic trigger and the condition $E\implies\property{}$; under a sufficiently small period, the condition $\property{}$ is evaluated exactly when $E$ holds.

\paragraph{Inspection objects and observations.}
The performance condition, $\property{}$ in \Cref{subsec:formal-guardrails}, operates over a finite sequence of system states, decisions, and/or events.
\lang exposes this behavior through a small set of typed \emph{inspection domains} (see Inspection rows of \Cref{tab:gdl-primitives}).
These include kernel objects (\textsf{KObject}), such as pages or processes; sets of kernel objects ($Set[\textsf{KObject}]$), which allow a guardrail to reason about collections as first-class entities (e.g., whether a property holds for all memory regions of 256 contiguous pages); and instrumented kernel functions (\textsf{KFunc}), which represent decision points or state transformations.
The choice of these domains is deliberate: guardrails targeted by \lang ultimately reason about kernel objects, their collections, and the function invocations that manipulate them.
By centering the language around these domains, \lang remains narrow while still covering a broad space of runtime performance conditions.

Over these domains, operators define \emph{observations} that extract the relevant runtime information.
This includes attributes of \textsf{KObject}s (e.g., whether a page has its {\tt dirty} bit set), set-level attributes of $Set[\ldots]$s (e.g., number of objects in the set), and decisions and execution times of instrumented functions.
The {\tt CompactionOverhead} example above illustrates the use of {\sf exec\_time}$(e)$ to extract the execution latency of {\tt compact\_zone()} invocations, while the {\tt BenefitCorrection} guardrail from \Cref{subsec:cbmm-case-study} illustrates reasoning over a $Set[\textsf{Page}]$ domain using the {\tt accesses} attribute and a {\sf read} from the in-kernel benefit map.
Any custom metric not representable through \textsf{KObject}s or \textsf{KFunc}s can be exposed via the {\sf read}(\textit{feat}) operation, which reads a pre-populated entry from an in-kernel store.
Note that guardrails are not responsible for collecting these features; the user must implement and maintain the relevant collection mechanisms.

\lang also supports aggregation ({\sf avg}, {\sf min}, {\sf max}, {\sf sum}) of observations and quantification ($\forall$, $\exists$) over elements of a domain.
This separation between inspection domains and observations is important: the former determines \emph{what objects} a guardrail looks at, while the latter determines \emph{what it measures} about them.
Operators write the performance condition as predicates over these observations.
Thus, although the semantics of guardrails are defined over execution traces, \lang does not expose raw traces directly; instead, it provides a structured interface for expressing performance conditions concisely.

\paragraph{Actions.}
A guardrail in \lang must also determine how the system should respond when a property is violated.
While the space of possible corrective actions is broad---different policies and different classes of violations may require different interventions---\lang intentionally exposes a narrow action interface supporting two types.
The first is a non-intrusive action, {\sf Report}, which records whenever a violation occurs for offline diagnosis; such reports can help in designing new policies altogether \ds{Add forward pointer to an example}.
The second is a corrective action, {\sf Set}$(flag, val)$, which modifies a control flag or parameter to alter subsequent system behavior.
Such updates may tune parameters, switch operating modes, select a different policy, or revert to a default configuration.

Both examples above illustrate the {\sf Set} action: {\tt CompactionOverhead} sets {\tt thp\_policy} to {\tt THPNever} when compaction is too expensive, and {\tt InDistributionShort} redirects to {\tt FLUXLongPredictor} when input features are out of distribution for the short-flow model.
Rather than encoding arbitrary recovery logic in the DSL, \lang uses these actions as a narrow interface to policy-specific mechanisms already implemented in the system.

\todo{Explain how the {\tt for} keyword defines the set of policies to which a guardrail applies.}
